\chapter{Discussion and Conclusion}
\label{ch:conclusion}

\section{Synthesis}
\label{sec:synth}

\paragraph{Hypothesis 1.} The claim was that financial infrastructure raises access more than usage
--- that expansion is formal rather than functional. Measured physically, as branches and ATMs, it
is \textbf{not supported}, and the more consequential result is why. Only 7.0 per cent of the
variation in branches per lakh adults, and 8.7 per cent of that in ATMs, occurs within states;
roughly 93 per cent is cross-sectional, which is precisely what fixed effects discard. The Hausman
test then rejects random effects, closing off the specification that would have used that
variation, and two defective observations out of 198 prove enough to reverse the sign of every ATM
coefficient. Doubling the window to twelve years doubles the usable variation and removes the
fragility, but leaves the verdict unchanged.

Measured digitally, the picture reverses. UPI adoption varies 67.1 per cent within states over the
same window, explanatory power rises roughly ninefold, and the predicted pattern appears:
$\Access{} +0.439$, $\Usage{} +0.183$, gap $+0.256$ at $p = 0.020$. The effect runs entirely
through credit accounts and not at all through deposit accounts --- the opposite of what an
accounting artefact would produce, since deposit accounts are the component entangled with UPI
registration by construction. The asymmetry is the finding. But UPI entered the design after the
physical null was seen, so this is a hypothesis generated by the data rather than one tested and
passed.

\stub{Summarise the verdicts on Hypotheses 2 and 3 here in the same form --- what was claimed, what
was found, what the result does not license --- then add a short paragraph on whether the three
point in a common direction. If they share the panel of Chapter~\ref{ch:data}, note explicitly
whether the within-state variation problem of Section~\ref{sec:h1-gate} affects them too: it is a
property of the data, not of Hypothesis~1.}

\section{Implications for policy}
\label{sec:policy}

\paragraph{Measure usage, not ownership.} Account ownership is near saturation, and the marginal
account opened now is very likely to be dormant: roughly 35 per cent of Indian adults with an
account hold an inactive one \citep{findex2021}, and about 26 per cent of PMJDY accounts are
reported inoperative \citep{pmjdy2025}. The RBI's Financial Inclusion Index already weights usage
(45 per cent) above access (35 per cent) \citep{rbi_fiindex2025}, and the evidence here supports
that ordering. Targets expressed in accounts opened measure the dimension that is no longer
binding.

\paragraph{The binding infrastructure has changed, and measurement should follow.} Branches per
lakh adults move a median of 5.6 per cent within a state over six years; UPI adoption moved by
orders of magnitude over the same period. Whatever the long-run effect of physical premises, they
are not the margin on which annual variation in financial inclusion was determined over
FY2019--FY2023. A framework that monitors branch and ATM density while treating digital payment
adoption as background is measuring the stationary part of the system.

\paragraph{If the credit channel is real, it is a policy object.} The digital result is
correlational, but the association it identifies --- digital payment adoption moving \emph{loan}
accounts while leaving deposit accounts untouched --- is consistent with payment history entering
credit underwriting, which is a deliberate feature of Indian financial architecture over this
period. That makes it a testable and governable channel rather than an incidental one, and it
sharpens what an evaluation should look for.

\paragraph{Data quality is itself a policy variable.} A discontinuity in one column of one RBI
return, affecting three states, was enough to reverse the sign of every ATM coefficient in Track~1
and to inflate a null deposit-account coefficient tenfold in Track~3. Series used for programme
evaluation warrant the revision and documentation discipline applied to headline aggregates.

\section{Implications for method}
\label{sec:method-implications}

\paragraph{Diagnose the identifying variation before estimating.} The variance decomposition was
specified as a gate, with thresholds fixed before the numbers were computed, and it correctly
forecast the instability the robustness checks later demonstrated. Had coefficients been examined
first, the null on \Access{} would have read as a substantive finding about infrastructure, and the
significant ATM coefficient under trimming as a discovery. Both readings would have been wrong, and
neither obviously so. The same statistic then did positive work: it is what identified digital
adoption as a regressor capable of answering the question, rather than merely another variable to
try. A wide confidence interval is ambiguous between ``there is no effect'' and ``there is no
variation from which to detect one'', and this distinction carries the whole report.

\paragraph{Decompose a composite outcome when any component is entangled with the regressor.} The
headline \Access{} coefficient on UPI would have been reportable as it stood. Splitting it revealed
that the entangled component was null and the clean component carried everything --- so the
composite was averaging a real effect with a defect-driven artefact, and the naive concern about
mechanical correlation pointed the wrong way. The rule for reading that decomposition was fixed
before it was run, which is what made the outcome interpretable either way.

\paragraph{Pre-registration is most valuable when results are ambiguous.} Had the estimates been
clean, the rule fixed in advance would have been redundant. They were not. It is also what forces
the honest statement of the digital result's status: the rule named branches and ATMs, so no
reading of the UPI finding can claim the hypothesis was tested and passed.

\paragraph{Slow-moving regressors and short panels are a recurring trap.} Infrastructure,
institutions and capital stocks all move slowly; two-way fixed effects is the default estimator;
and the Hausman test will typically forbid retreat to random effects in exactly the settings where
between-state variation is most correlated with the outcome. The reconsideration of
\citet{burgess2005} by \citet{buliskeria2022} makes the same point from the other direction. The
remedy applied here --- find the infrastructure that actually moved --- is more generally available
than it looks.

\section{Directions for further work}
\label{sec:further}

Three things would strengthen this, in order of value. \textbf{An identification strategy for UPI}
is the first and largest: the current digital result is a within-state correlation between two
co-trending series, and candidate instruments include state-level 4G rollout timing or the
staggered geography of Jio's 2016--17 entry. \textbf{Total UPI rather than one firm} would remove
the varying-market-share problem, since NPCI publishes state-level aggregates that would replace
the PhonePe proxy. \textbf{Extending the digital panel} would take it from five years to eight, the
PhonePe series running to FY2026 while the banking series stop at FY2023. Beyond these,
sub-state disaggregation would raise precision substantially, and three measurement problems are
tractable: replacing the interpolated adult denominator with Census-based age shares, estimating on
logs or ranks rather than pooled min--max levels to address the compression of \Usage{}, and
reconciling the three attribution rules currently merged into one panel.

\section{Conclusion}
\label{sec:conclusion}

This report asked whether the expansion of financial infrastructure in India has been formal rather
than functional --- whether it opens accounts faster than it puts them to use. The question is well
motivated: India has effectively solved account ownership while recording the highest rate of
account dormancy in the world.

Asked of \emph{physical} infrastructure, the question cannot be answered with these data.
Hypothesis~1 is not supported, but the null is uninformative: branches and ATMs do not move enough
within Indian states over six years, or even twelve, for a fixed-effects panel to detect their
effect, and the specification test that would license using cross-sectional variation instead
rejects that estimator decisively. Two observations out of 198 reverse the sign of every ATM
coefficient, which demonstrates the point rather than merely asserting it.

Substituting the infrastructure that actually moved changes the picture. UPI adoption varies
abundantly within states, and on that measure the predicted pattern appears: access rises more than
usage, and the gap between them is significantly different from zero for the first time in the
study. It does so entirely through credit accounts rather than deposit accounts --- an asymmetry
that is the reverse of what a mechanical artefact would generate, and one that fits a documented
channel, digital payment history feeding credit underwriting.

\begin{keybox}{The honest headline}
A state panel cannot identify the effect of \emph{physical} banking infrastructure on financial
inclusion, because branches and ATMs barely move within states. Substituting the infrastructure
that actually moved --- digital payments --- the picture reverses: UPI adoption predicts access
more strongly than usage, significantly so, and it does it entirely through credit accounts rather
than deposit accounts. \textbf{The identification problem is solved; the causal one is not.}
\end{keybox}

\noindent
Two boundaries must travel with that. The pre-committed hypothesis named branches and ATMs and it
failed; UPI was added afterwards, so the digital finding is a hypothesis this data \emph{generated}
and not one it tested. And fixed effects remove fixed state differences and the common national
trend, not state-specific time-varying confounders --- smartphone penetration, fintech entry and
youth share all rose alongside both UPI adoption and credit accounts.

What remains is a clear direction rather than a closed question. The variation this design needed
and did not have exists in abundance on the digital side; what it now lacks is exogeneity, not
power. That is a more tractable problem than the one it started with, and the discipline that made
this report's negative result legible --- fixing the decision rule and the diagnostic thresholds
before seeing the numbers --- is what should be carried into solving it.
